\section{Master in a University at War with Its Friars}

\subsection{The delayed chair}

Thomas qualified for the mastership in theology in 1256, younger than the
ordinary statutory age.  His advancement occurred amid a bitter conflict
between secular masters and the mendicant orders.  The dispute involved
teaching posts, corporate oaths, pastoral privileges, poverty, and the right of
religious institutes to participate in the university.  William of Saint-Amour
attacked the friars not merely as professional competitors but as dangers to
the Church.  Papal intervention by Alexander~IV and the condemnation of
William's \emph{De periculis novissimorum temporum} helped secure the mendicant
masters' position.

Thomas's \emph{Against Those Who Attack the Worship of God and Religion}
answers within that conflict.  He defends a religious life ordered to study,
teaching, and preaching; voluntary poverty does not make apostolic ministry
illicit, and contemplation can overflow into instruction.  The polemical
occasion matters.  A later reader should not turn a defense of new orders into
a timeless denial of institutional self-interest, nor reduce the theological
argument to a fight over jobs.  Thomas believed the mendicant form could serve
the Church because disciplined study and common poverty freed preaching for
truth.

The received inaugural lecture associated with the verse ``watering the
mountains from above'' presents sacred teachers as receiving wisdom from God
and communicating it without owning its source.  Whether every detail of the
occasion is fixed or not, the image fits the office he practiced: authority is
ministerial, ordered downward from divine truth through finite teachers to
hearers.  It is a demanding ideal for someone whose name would later function
as an authority in its own right.

\subsection{The first Paris regency}

From 1256 to 1259 Thomas taught as master, disputed publicly, preached, and
supervised the academic exercises attached to his chair.  Works commonly
placed here include disputed questions on truth, the beginning of substantial
biblical exposition, and defenses of religious life.  The university method
did not mean assembling token objections before a predetermined slogan.  A
good article identifies the issue, states serious reasons on both sides,
distinguishes senses, gives a determination, and returns to each objection.
The form trained intellectual charity when practiced honestly: an opponent's
reason had to become intelligible before it could be answered.

The same form could sharpen combat.  Thomas wrote within real institutional
struggle and could describe adversaries through categories they rejected.  His
later reputation for serenity must not bleach the argumentative stakes.  Yet
his characteristic answer was architectural rather than merely victorious.
Natural reason, revealed teaching, vowed perfection, ecclesial office, and
pastoral usefulness had to be placed in an order in which one good did not
destroy another.

In 1259 Thomas left the Paris chair and returned to Italy under Dominican
obedience.  The move ended neither authorship nor public responsibility.  It
shifted their geography and brought him closer to the papal court, Greek-Latin
ecclesial questions, Dominican formation, and the liturgical work with which
later memory most closely joined his prayer and poetry.
